\graphicspath{{content/chapters/4_evaluation/figures/}}

\chapter{Evaluation}
\label{chp:evaluation}

This chapter presents the empirical results of the analyses 
described in the methodology. The evaluation proceeds in four 
stages. First, the demographic distributions of the annotated 
Re-LAION-5B and Stable Diffusion v1.5 corpora are 
characterised at both aggregate and profession level. Second, 
the dataset classification results from the spurious feature 
transformation pipeline are reported and interpreted, with 
per-class accuracy attribution used to localise the source 
of discriminative signal within each transformation. Third, 
the demographic probing results are presented, assessing 
whether gender and skin tone can be inferred from low-level 
image representations independently of semantic content. 
Fourth, the debiased generation corpus is evaluated against 
the aggregate Stable Diffusion v1.5 results to assess 
whether the inference-time interventions produce measurable 
demographic shifts. The chapter concludes with a unified 
discussion of findings, their implications, and the 
limitations of the evaluation.

\section{Demographic Distribution Analysis}
\label{sec:demographic_distributions}

\subsection{Aggregate Gender and Skin Tone Distributions}

Figures~\ref{fig:gender_dist_aggregate} and 
\ref{fig:skintone_dist_aggregate} show the aggregate gender 
and skin tone distributions across all 200 professions for 
both Re-LAION-5B and Stable Diffusion v1.5. 
Figure~\ref{fig:mst_dist_aggregate} shows the full MST 1--10 
skin tone distribution for both datasets.

\begin{figure}[htbp]
	\centering
	\includegraphics[width=0.7\textwidth]{aggregate_gender.png}
	\caption{Aggregate gender distribution across all 200 
		professions for Re-LAION-5B and Stable Diffusion v1.5. 
		The dashed line indicates demographic parity at 50\%.}
	\label{fig:gender_dist_aggregate}
\end{figure}

\begin{figure}[htbp]
	\centering
	\includegraphics[width=\textwidth]{aggregate_skintone_bins.png}
	\caption{Aggregate skin tone distribution using the 
		Light (MST 1--3), Mid (MST 4--7), and Dark (MST 8--10) 
		three-bin grouping for Re-LAION-5B and Stable Diffusion 
		v1.5.}
	\label{fig:skintone_dist_aggregate}
\end{figure}

\begin{figure}[htbp]
	\centering
	\includegraphics[width=\textwidth]{aggregate_mst_distribution.png}
	\caption{Aggregate MST 1--10 skin tone distribution for 
		Re-LAION-5B and Stable Diffusion v1.5. Shaded regions 
		indicate the Light (MST 1--3), Mid (MST 4--7), and Dark 
		(MST 8--10) bin boundaries.}
	\label{fig:mst_dist_aggregate}
\end{figure}

Both datasets exhibit a pronounced male skew across the full 
200-profession corpus, but the imbalance is substantially more 
severe in Stable Diffusion v1.5 than in Re-LAION-5B. 
Re-LAION-5B contains 83,003 female-labelled images (41.5\%) 
against 116,997 male-labelled images (58.5\%), a gap of 17 
percentage points relative to parity. Stable Diffusion v1.5 
outputs widen this gap to 42 percentage points, with only 
58,037 female-labelled images (29.02\%) against 141,963 
male-labelled images (70.98\%). The generative model therefore 
does not merely reflect the gender distribution of its training 
corpus but amplifies it substantially, producing outputs that 
are considerably more male-skewed than the web-scraped imagery 
from which it learned.

The skin tone distributions reveal a structurally different 
pattern. Both datasets are overwhelmingly dominated by 
Mid-bin predictions (MST 4--7), with Re-LAION-5B at 86.74\% 
and Stable Diffusion v1.5 at 76.99\%. The MST 1--10 
distribution in Figure~\ref{fig:mst_dist_aggregate} reveals 
that within the Mid bin, the two datasets differ in their 
modal value: Re-LAION-5B peaks at MST 5 and MST 6, while Stable Diffusion v1.5 peaks 
at MST 4 and MST 5, representing a shift toward lighter 
tones within the Mid range. The more pronounced divergence 
lies at the extremes: Re-LAION-5B has a Light-bin proportion 
of 12.45\% and a Dark-bin proportion of 0.81\%, indicating 
that non-Mid predictions are distributed, albeit very 
unevenly, across both ends of the tone scale. Stable Diffusion 
v1.5 shows a substantially higher Light-bin proportion of 
22.87\% but a Dark-bin proportion of just 0.14\% --- a 
reduction of 83\% relative to Re-LAION-5B. MST values 9 and 
10 are effectively absent in both datasets, with Stable 
Diffusion v1.5 generating near-zero proportions even at 
MST 8. This near-complete collapse of dark skin tone 
representation in the generative outputs, combined with the 
shift toward lighter tones within the Mid bin, suggests that 
the model has absorbed a strong light skin tone prior from 
its training corpus and concentrates its outputs toward 
the lighter half of the tonal range.

\subsection{Per-Profession and Per-ISCO-Group Analysis}

Figures~\ref{fig:gender_by_isco},~\ref{fig:skintone_by_isco}, 
and~\ref{fig:mst_heatmap_isco} break down the gender and skin 
tone distributions by ISCO-08 major occupational group. 
Tables~\ref{tab:top_biased_professions_gender} 
and~\ref{tab:top_biased_professions_skin} list the ten 
professions with the strongest observed gender and light skin 
tone skew in each dataset respectively.

\begin{figure}[htbp]
	\centering
	\includegraphics[width=\textwidth]{isco_gender.png}
	\caption{Female proportion by ISCO-08 major occupational 
		group for Re-LAION-5B and Stable Diffusion v1.5. 
		The dashed line indicates demographic parity at 50\%.}
	\label{fig:gender_by_isco}
\end{figure}

\begin{figure}[htbp]
	\centering
	\includegraphics[width=\textwidth]{isco_skintone_bins.png}
	\caption{Skin tone distribution (Light, Mid, Dark bins) 
		by ISCO-08 major occupational group for Re-LAION-5B 
		and Stable Diffusion v1.5.}
	\label{fig:skintone_by_isco}
\end{figure}

\begin{figure}[htbp]
	\centering
	\includegraphics[width=\textwidth]{isco_mst_heatmap.png}
	\caption{MST 1--10 score distribution (\%) by ISCO-08 
		major occupational group for Re-LAION-5B and Stable 
		Diffusion v1.5. Cell values indicate the proportion 
		of images in each group assigned to each MST score.}
	\label{fig:mst_heatmap_isco}
\end{figure}

\begin{table}[htbp]
	\centering
	\caption{Ten professions with the strongest gender skew 
		in Re-LAION-5B and Stable Diffusion v1.5. Gender 
		skew is the absolute deviation of the female proportion 
		from 0.5; higher values indicate stronger skew 
		regardless of direction. F\% = female percentage; 
		M\% = male percentage.}
	\label{tab:top_biased_professions_gender}
	\small
	\begin{tabular}{llrrr}
		\toprule
		\textbf{Dataset} & \textbf{Profession} & 
		\textbf{F\%} & \textbf{M\%} & \textbf{Skew} \\
		\midrule
		Re-LAION-5B & Administrative Assistant & 94.7 & 5.3  & 0.447 \\
		Re-LAION-5B & Makeup Artist            & 92.7 & 7.3  & 0.427 \\
		Re-LAION-5B & Receptionist             & 92.0 & 8.0  & 0.420 \\
		Re-LAION-5B & Executive Assistant      & 91.5 & 8.5  & 0.415 \\
		Re-LAION-5B & Office Assistant         & 90.2 & 9.8  & 0.402 \\
		Re-LAION-5B & Dietitian                & 88.6 & 11.4 & 0.386 \\
		Re-LAION-5B & Paralegal                & 87.8 & 12.2 & 0.378 \\
		Re-LAION-5B & Bodyguard                & 13.0 & 87.0 & 0.370 \\
		Re-LAION-5B & Priest                   & 13.3 & 86.7 & 0.367 \\
		Re-LAION-5B & Blacksmith               & 14.5 & 85.5 & 0.355 \\
		\midrule
		SDv1.5 & Logger                        &  2.5 & 97.5 & 0.475 \\
		SDv1.5 & Nurse                         & 97.5 &  2.5 & 0.475 \\
		SDv1.5 & Receptionist                  & 97.3 &  2.7 & 0.473 \\
		SDv1.5 & Farmer                        &  3.0 & 97.0 & 0.470 \\
		SDv1.5 & Dietitian                     & 96.9 &  3.1 & 0.469 \\
		SDv1.5 & Makeup Artist                 & 96.9 &  3.1 & 0.469 \\
		SDv1.5 & Flight Attendant              & 96.6 &  3.4 & 0.466 \\
		SDv1.5 & Administrative Assistant      & 96.4 &  3.6 & 0.464 \\
		SDv1.5 & Construction Worker           &  3.6 & 96.4 & 0.464 \\
		SDv1.5 & Miner                         &  3.9 & 96.1 & 0.461 \\
		\bottomrule
	\end{tabular}
\end{table}

\begin{table}[htbp]
	\centering
	\caption{Ten professions with the highest Light-bin (MST 
		1--3) skin tone proportion in Re-LAION-5B and Stable 
		Diffusion v1.5.}
	\label{tab:top_biased_professions_skin}
	\small
	\begin{tabular}{llrrr}
		\toprule
		\textbf{Dataset} & \textbf{Profession} & 
		\textbf{Light\%} & \textbf{Mid\%} & \textbf{Dark\%} \\
		\midrule
		Re-LAION-5B & Makeup Artist              & 27.6 & 72.3 & 0.1 \\
		Re-LAION-5B & Hair Dresser               & 21.8 & 78.0 & 0.2 \\
		Re-LAION-5B & Dentist                    & 20.4 & 79.3 & 0.3 \\
		Re-LAION-5B & Carer                      & 19.9 & 79.8 & 0.3 \\
		Re-LAION-5B & Florist                    & 18.7 & 81.1 & 0.2 \\
		Re-LAION-5B & Singer                     & 18.6 & 81.1 & 0.3 \\
		Re-LAION-5B & Barber                     & 18.5 & 80.2 & 1.3 \\
		Re-LAION-5B & Customer Support Specialist & 18.5 & 81.2 & 0.3 \\
		Re-LAION-5B & Comedian                   & 18.1 & 81.0 & 0.9 \\
		Re-LAION-5B & Actor                      & 17.2 & 80.7 & 2.1 \\
		\midrule
		SDv1.5 & Librarian             & 56.8 & 43.1 & 0.1 \\
		SDv1.5 & Carer                 & 55.5 & 44.4 & 0.1 \\
		SDv1.5 & Politician            & 50.2 & 49.8 & 0.0 \\
		SDv1.5 & Meteorologist         & 45.3 & 54.7 & 0.0 \\
		SDv1.5 & Mayor                 & 45.2 & 53.9 & 0.9 \\
		SDv1.5 & News Reader           & 45.2 & 54.7 & 0.1 \\
		SDv1.5 & Farmer                & 43.7 & 56.3 & 0.0 \\
		SDv1.5 & Comedian              & 41.3 & 58.7 & 0.0 \\
		SDv1.5 & Florist               & 40.5 & 59.5 & 0.0 \\
		SDv1.5 & Conservation Officer  & 40.4 & 59.1 & 0.5 \\
		\bottomrule
	\end{tabular}
\end{table}

The ISCO-08 group breakdown in Figure~\ref{fig:gender_by_isco} 
reveals that gender skew is pervasive across the majority of occupational 
categories across both datasets, but its magnitude and pattern 
differ markedly. In Re-LAION-5B, Clerical Support is the only 
major group to exceed demographic parity at 58.26\% female, 
consistent with the historically female-dominated composition 
of administrative and secretarial roles. All remaining groups 
fall below parity, with Plant \& Machine Operators showing 
the strongest male skew at 77.85\% male, followed by Armed 
Forces at 76.5\% and Managers at 64.53\%. In Stable Diffusion 
v1.5, the direction of skew is broadly preserved --- Clerical 
Support remains the only group at or near parity at 51.93\% 
female --- but the magnitude increases substantially across 
every group without exception. The most severe amplifications 
occur in Skilled Agricultural (from 43.3\% to 7.3\% female, 
a drop of 36 percentage points), Armed Forces (from 23.5\% 
to 7.7\%, a drop of 15.8 percentage points), and Plant \& 
Machine Operators (from 22.15\% to 13.55\%). This systematic 
reduction in female representation across all occupational 
categories confirms that the amplification effect identified 
at the aggregate level is not driven by a small number of 
outlier professions but reflects a broad structural tendency 
of the generative model.

Comparing these distributions against the US BLS 2024 and
Global ILOSTAT reference markers in
Figure~\ref{fig:gender_by_isco} reveals a consistent pattern
of under-representation of women relative to real-world labour
force composition. Both reference baselines show broad
agreement in their directional signal: groups identified as female‐majority in US
BLS data are also female‐majority in the global statistics, while groups such as Skilled
Agricultural, Craft \& Trades, and Plant \& Machine Operators are male‐dominated in
both.. This convergence between the two independent baselines strengthens the
validity of the comparison and indicates that the observed
gaps are not artefacts of geographic or methodological
differences between the two reference sources.

Against these baselines, Re-LAION-5B presents a mixed picture that depends on the reference used. Relative to US BLS 2024, Re-LAION-5B falls below the real-world female proportion in approximately six of ten ISCO groups, most notably in Professionals (approximately 43\% vs. 50\%), Service \& Sales (approximately 46\% vs. 51\%), and Technicians \& Associate Professionals (approximately 44\% vs. 49\%). Against the Global ILOSTAT baseline, the picture is considerably more balanced, with Re-LAION-5B exceeding the global female proportion in several groups including Skilled Agricultural, Craft \& Trades, and Plant \& Machine Operators, while falling below it in Professionals and Technicians \& Associate Professionals. This divergence between the two baselines reflects genuine differences in occupational gender composition between the US and global labour markets rather than inconsistency in the dataset, and suggests that Re-LAION-5B's gender distributions are broadly consistent with global workforce composition in some occupational domains while remaining below US-specific norms in others.

Taken together, these comparisons establish a two-stage bias pattern. Re-LAION-5B exhibits gender distributions that diverge from real-world labour force composition in a baseline-dependent manner, falling below US-specific norms in the majority of occupational groups while remaining broadly consistent with global workforce proportions in several others. Stable Diffusion v1.5 then amplifies the male skew further across nearly every group regardless of baseline, producing occupational gender distributions that are more male-skewed than even the already imbalanced training corpus. The generative model therefore does not merely inherit the distributional properties of its training data but actively exacerbates them, producing representations that diverge substantially from both demographic parity and empirically observed occupational gender composition.

At the individual profession level, 167 of 200 professions 
(83.5\%) show gender skew in the same direction in both 
datasets, confirming that Stable Diffusion v1.5 largely 
inherits the directional gender associations of its training 
corpus. However, in 178 of 200 professions (89.0\%) the 
magnitude of skew is greater in Stable Diffusion v1.5 than 
in Re-LAION-5B. The most skewed professions in 
Table~\ref{tab:top_biased_professions_gender} illustrate this 
pattern clearly: Administrative Assistant, which already shows 
a strong female skew in Re-LAION-5B at 94.7\%, reaches 96.4\% 
in SDv1.5; Dietitian rises from 88.6\% to 96.9\%; and 
Receptionist rises from 92.0\% to 97.3\%. At the male-skewed 
end, the stereotyping is equally pronounced, with Logger 
(97.5\% male), Farmer (97.0\%), and Construction Worker 
(96.4\%) among the most extreme SDv1.5 outputs. The 33 professions where the datasets disagree in skew direction most likely reflect cases where the training corpus signal is weak or near parity, such that small generative shifts produce a nominal change in direction. These cases warrant further investigation but are beyond the scope of the current analysis.

The skin tone analysis at the ISCO group level, shown in 
Figures~\ref{fig:skintone_by_isco} and~\ref{fig:mst_heatmap_isco}, 
presents a consistent story across all occupational categories. 
In Re-LAION-5B the Light-bin proportion varies narrowly 
between 6.9\% (Armed Forces) and 13.59\% (Clerical Support), 
indicating that no occupational group is strongly associated 
with light skin tones in the training corpus. In Stable 
Diffusion v1.5, all ISCO groups show a substantially 
elevated Light-bin proportion, ranging from 16.85\% (Plant \& 
Machine Operators) to 35.1\% (Skilled Agricultural). The 
MST heatmap confirms that this shift is not caused by a 
redistribution within the Light bin but by a leftward shift 
in the modal MST score: while Re-LAION-5B concentrates mass 
at MST 5--6 across all groups, Stable Diffusion v1.5 
systematically concentrates mass at MST 4--5, with 
correspondingly elevated MST 3 proportions. The Dark-bin 
collapse observed at the aggregate level is equally 
consistent across ISCO groups: Re-LAION-5B dark proportions 
range from 0.5\% to 1.8\%, while in SDv1.5 they approach 
zero in most groups, with Skilled Agricultural recording 0.0\% 
dark-bin images entirely.

At the individual profession level, the scale of the skin tone 
amplification in SDv1.5 is particularly stark. The most 
light-skewed profession in Re-LAION-5B is Makeup Artist at 
27.6\% Light, yet in SDv1.5 this figure rises to 55.5\% for 
Carer and 56.8\% for Librarian --- professions that in 
Re-LAION-5B reach only 19.9\% and are absent from the top 10 
entirely. The SDv1.5 top 10 light-skewed professions all 
exceed 40\% Light, a threshold not reached by any profession 
in Re-LAION-5B. In 175 of 200 professions (87.5\%), SDv1.5 
generates a higher proportion of Light-bin images than 
Re-LAION-5B does for the same profession, establishing a 
near-universal tendency toward lighter skin tone generation 
that is independent of occupational category.

Unlike the gender analysis, a direct comparison of skin tone
distributions against real-world reference statistics was not
conducted. The BLS CPS data reports labour force composition
in terms of self-identified racial categories rather than
skin tone measurements, and converting race proportions into
expected MST distributions requires an intermediate
race-to-skin-tone mapping. As described in
Section~\ref{subsec:race_to_mst}, this mapping was derived
empirically from the FairFace dataset using the VGG-16 MST
regression model; however, inspection of the resulting
conditional distributions revealed that all racial groups map
predominantly to the Mid bin (MST 4--7), reflecting the tonal
range of the FairFace corpus rather than the full MST scale.
The derived per-profession reference distributions therefore
lack sufficient discriminative resolution to serve as
meaningful baselines for the skin tone comparison, collapsing
to near-identical Mid-dominant profiles regardless of
occupational racial composition. As a result, skin tone
analysis in this work is evaluated solely against the
demographic parity baseline. The construction of a
more granular and reliable race-to-skin-tone reference
distribution, whether through a larger and more tonally
diverse calibration dataset or through direct MST annotations
on labour force samples, remains an open challenge and
constitutes a limitation of the present analysis.

%% ============================================================
%% Section 2: Dataset Classification Results
%% ============================================================
\section{Dataset Classification}

Dataset classification experiments assess whether a ConvNeXt-Tiny 
classifier trained on transformed image representations can 
distinguish Re-LAION-5B, Stable Diffusion v1.5, and COCO images 
as a three-class problem. Each transformation is treated as an 
independent experiment: images are transformed, split into training, 
validation, and test sets, and the classifier is trained on the 
transformed training set with early stopping (patience = 2) and 
evaluated on the held-out test set. Chance level for all experiments 
is $\nicefrac{1}{3} \approx 33.3\%$ under a balanced three-class 
prior. Above-chance accuracy on a given transformation indicates that 
the transformation retains dataset-discriminating signal; the 
magnitude of accuracy quantifies how strongly that signal is 
preserved. Classification accuracy is used as the primary metric 
throughout this section, consistent with the three-class evaluation 
protocol of Zeng et al. \cite{Zeng_2024_UnderstandingBiasinLargeScaleVisualDatasets} from 
which the transformation framework is adopted. 

COCO is included solely to extend the classification task beyond a binary comparison 
between Re-LAION-5B and Stable Diffusion v1.5, enabling per-class accuracy attribution 
that identifies which dataset drives separability under each transformation. As COCO 
carries no profession-conditioned annotation for gender or skin tone within this 
pipeline, it does not appear in the demographic analyses that follow; the results in 
this section therefore reflect general distributional separability across datasets and 
should not be interpreted as direct evidence of demographic bias, which is examined 
independently in Section~\ref{sec:demographic_attribute_separability}.

Table~\ref{tab:dataset_classification_results} summarises the results for all transformations evaluated. While COCO was included to enable per-class attribution, explicit per-class breakdowns are not reported given the uniformly near-perfect AUC values across virtually all transformations — the majority exceeding 0.99 — making it highly unlikely that any individual class is separated substantially below the overall accuracy level. The aggregate metrics are therefore considered sufficient to characterise the separability of the full three-class problem.

\begin{table}[htbp]
	\centering
	\caption{Dataset classification results per image transformation (three-class:
		Re-LAION-5B, Stable Diffusion v1.5, COCO). All models trained with ConvNeXt-Tiny,
		early stopping (patience = 2), seed = 0.}
	\label{tab:dataset_classification_results}
	\small
	\begin{threeparttable}
	\begin{tabular}{p{4cm}cccc}
		\toprule
		\textbf{Transformation} &
		\textbf{Val Acc (\%)} & \textbf{Val AUC} &
		\textbf{Test Acc (\%)} & \textbf{Test AUC} \\
		\midrule
		Base (unmodified)          & 99.56 & 0.9996 & 99.42 & 0.9996 \\
		VAE reconstruction         & 99.15 & 0.9992 & 99.15 & 0.9992 \\
		High-pass filter           & 98.65 & 0.9991 & 98.64 & 0.9988 \\
		Patch shuffle $8\times8$   & 98.41 & 0.9989 & 97.91 & 0.9983 \\
		Patch shuffle $16\times16$ & 98.60 & 0.9986 & 97.69 & 0.9983 \\
		Low-pass filter            & 97.30 & 0.9977 & 96.99 & 0.9969 \\
		Depth                      & 96.41 & 0.9968 & 96.78 & 0.9965 \\
		Patch shuffle $4\times4$   & 97.20 & 0.9968 & 94.38 & 0.9949 \\
		Edge                       & 95.17 & 0.9934 & 93.62 & 0.9905 \\
		Semantic segmentation      & 90.47 & 0.9814 & 91.36 & 0.9837 \\
		Captions$^\dagger$         & ---   & ---    & 90.99 & 0.9820 \\
		Patch shuffle $2\times2$   & 88.97 & 0.9839 & 89.91 & 0.9766 \\
		Object detection           & 85.88 & 0.9617 & 88.62 & 0.9728 \\
		Pixel shuffle              & 86.01 & 0.9574 & 80.67 & 0.9430 \\
		Mean RGB                   & 70.80 & 0.8443 & 64.19 & 0.8207 \\
		\bottomrule
	\end{tabular}
	\begin{tablenotes}
		\footnotesize
		\item[$\dagger$] Logistic regression probe on BLIP Large captions encoded via
		SentenceT5-Base; all other rows use ConvNeXt-Tiny on rendered images.
	\end{tablenotes}
	\end{threeparttable}

\end{table}

\paragraph{1. The three datasets are separable across every level of representation tested}

Every transformation achieves well above the chance level of 33.3\%, with even the most aggressive information compression — mean RGB at 64.19\% — retaining substantial discriminating signal. This demonstrates that dataset-level differences between Re-LAION-5B, Stable Diffusion v1.5, and COCO persist even under extreme reduction of image content. Even when spatial structure, semantics, and most visual detail are removed, the classifier can still attribute images to their source dataset significantly better than chance.

The persistence of separability across such compressed representations suggests that dataset identity is encoded in multiple correlated visual cues rather than a single isolated artefact. These cues may include colour statistics, spatial patterns, semantic composition, or other structural regularities that differ systematically between datasets. Because these signals remain detectable even after large portions of the visual information are removed, they likely reflect broad statistical regularities rather than isolated features.

Importantly, these experiments demonstrate separability rather than its causal origins. The observed differences may arise from a combination of dataset curation practices, object distribution differences, scene composition biases, or visual statistics introduced by generative models. Consequently, the results indicate that the datasets possess distinct statistical signatures without directly identifying the specific factors responsible for those differences.

\paragraph{2. Dataset separability cannot be explained by acquisition artefacts alone}

The VAE result (99.15\%) is virtually identical to the unmodified base (99.42\%), a reduction of only 0.27 percentage points. Variational autoencoder reconstruction compresses each image through a latent bottleneck before reconstructing it, which typically suppresses certain low-level acquisition artefacts such as compression noise or encoding irregularities while preserving perceptually meaningful image structure.

Because VAE reconstruction largely preserves colour distributions, texture statistics, and spatial structure while reducing some superficial artefacts associated with file encoding or compression, the near-identical classification performance suggests that simple acquisition artefacts are unlikely to be the dominant source of discriminating signal. If dataset separability were primarily driven by compression differences, file-format characteristics, or other encoding artefacts, the VAE reconstruction process would be expected to reduce classification accuracy more substantially.

Instead, the results suggest that dataset differences arise primarily from content and structural properties of the images themselves. These may include differences in scene composition, object arrangement, photographic style, or statistical regularities introduced during dataset creation. However, it should be noted that VAE reconstruction does not eliminate all low-level image statistics, and therefore the results indicate that acquisition artefacts are unlikely to dominate rather than definitively ruling out all low-level contributions.

\paragraph{3. Local spatial statistics are highly dataset-discriminative}

Even $2\times2$ patch shuffling — which destroys global layout while preserving local spatial structure including texture patterns, gradient relationships, and colour correlations within each patch — achieves 89.91\% test accuracy. Increasing patch size substantially improves performance, with $8\times8$ patch shuffling reaching 97.91\% and $16\times16$ achieving 97.69\%.

Patch shuffling disrupts global object arrangement and scene layout while preserving local spatial statistics within each patch. These statistics include short-range correlations between neighbouring pixels, micro-texture patterns, and small-scale colour gradients that often characterise the visual appearance of images. Because these local structures remain intact even when global spatial coherence is removed, classifiers trained on shuffled images must rely primarily on these fine-scale statistical cues.

The strong performance under patch shuffling therefore indicates that dataset-specific signatures are embedded in local spatial statistics rather than solely in high-level scene composition. Such statistics may arise from differences in photographic style, image processing pipelines, or generative model priors that produce consistent patterns at fine spatial scales. However, the experiment does not isolate the origin of these statistics, and the observed separability may reflect a combination of generative artefacts, dataset curation biases, or inherent distributional differences between the source corpora.

\paragraph{4. Both frequency bands carry strong discriminating signal, with a modest advantage for high-frequency components}

High-pass filtering achieves 98.64\% test accuracy, outperforming low-pass filtering at 96.99\% by 1.65 percentage points. High-frequency components correspond primarily to edges, fine textures, and small-scale spatial variations, whereas low-frequency components capture coarse illumination gradients and large-scale structural information.

The strong performance under both filtering operations indicates that dataset-level differences are distributed across multiple spatial scales. High-frequency signals such as edge patterns and texture statistics appear to carry slightly more discriminating information than low-frequency signals, suggesting that fine-grained visual details contribute strongly to dataset identity. At the same time, the high performance under low-pass filtering indicates that broader structural cues such as global shape patterns or scene layouts also play an important role.

The relatively small difference between the two results suggests that dataset bias is not concentrated exclusively in either frequency band but instead permeates the entire frequency spectrum.

\paragraph{5. A hierarchy of dataset-discriminating signal exists across representational levels}

% Table~\ref{tab:dataset_classification_results} reveals a consistent ordering of 
% discriminating power from strongest to weakest across the full transformation suite. 
% Unmodified and VAE-reconstructed images achieve near-perfect accuracy ($\geq$99\%), 
% confirming that the raw statistical properties of each dataset are highly distinctive. 
% Frequency and fine spatial statistics --- high-pass filter, $8{\times}8$ and 
% $16{\times}16$ patch shuffle --- follow closely at 97--99\%, indicating that 
% dataset-specific signatures are preserved in local spatial structure even when global 
% layout is destroyed. Structural geometry transformations --- depth and Canny edge --- 
% occupy the next tier at 93--97\%, demonstrating that scene geometry alone is 
% sufficient for reliable classification without any colour or texture information. 
% Dense semantic representations --- segmentation at 91.36\% --- and semantic content 
% probes --- captions at 90.99\% and object detection at 88.62\% --- follow, showing 
% that object identity and linguistic content carry substantial but somewhat weaker 
% discriminating power than structural cues. Global colour statistics form the lowest 
% tier: $2{\times}2$ patch shuffle at 89.91\%, pixel shuffle at 80.67\%, and mean RGB 
% at 64.19\%, the last of which retains only the three-dimensional average colour vector 
% yet still achieves 31 percentage points above chance. This hierarchy indicates that 
% dataset bias is distributed across multiple correlated visual attributes rather than 
% concentrated in any single representation, and that structural and low-level 
% statistical differences between the datasets are more pronounced than semantic 
% content differences. Critically, no single transformation eliminates separability, even the most aggressive compression retains substantial above-chance performance confirming that the statistical differences between datasets are not attributable to any single visual property but are instead redundantly encoded across colour, texture, spatial structure, geometry, semantics, and language simultaneously.

Table~\ref{tab:dataset_classification_results} reveals a consistent ordering of discriminating power from strongest to weakest across the full transformation suite. Unmodified and VAE-reconstructed images achieve near-perfect accuracy ($\geq$99\%), confirming that the raw statistical properties of each dataset are highly distinctive. Frequency and fine spatial statistics --- high-pass filter, $8{\times}8$ and $16{\times}16$ patch shuffle --- follow closely at 97--99\%, indicating that dataset-specific signatures are preserved in local spatial structure even when global layout is destroyed. Patch shuffle $4{\times}4$ at 94.38\% sits at the boundary between the spatial and geometry tiers, reflecting the transitional nature of this patch size between fine local statistics and broader structural information. Structural geometry transformations --- depth and Canny edge --- occupy the next tier at 93--97\%, demonstrating that scene geometry alone is sufficient for reliable classification without any colour or texture information. Dense semantic representations --- segmentation at 91.36\% --- and semantic content probes --- captions at 90.99\% and object detection at 88.62\% --- follow, showing that object identity and linguistic content carry substantial but somewhat weaker discriminating power than structural cues. Global colour statistics form the lowest tier: pixel shuffle at 80.67\% and mean RGB at 64.19\%, the last of which retains only the three-dimensional average colour vector yet still achieves 31 percentage points above chance. Patch shuffle $2{\times}2$ at 89.91\% sits at the boundary between the spatial and colour tiers, retaining minimal local structure but performing above the purely colour-based representations. This hierarchy indicates that dataset bias is distributed across multiple correlated visual attributes rather than concentrated in any single representation, and that structural and low-level statistical differences between the datasets are more pronounced than semantic content differences. Critically, no single transformation eliminates separability --- even the most aggressive compression retains substantial above-chance performance --- confirming that the statistical differences between datasets are not attributable to any single visual property but are instead redundantly encoded across colour, texture, spatial structure, geometry, semantics, and language simultaneously.

\paragraph{6. Semantic dataset differences survive compression to natural language}

Caption-based classification using SentenceT5 embeddings achieves 90.99\% test accuracy. Because the classifier receives only textual descriptions rather than visual input, this result indicates that systematic differences between the datasets are sufficiently consistent to be reflected in natural language descriptions of the images.

Differences in object co-occurrence patterns, scene composition, and contextual relationships between objects may therefore manifest in textual representations as well as in visual statistics. The persistence of dataset separability under this modality shift suggests that dataset identity is not solely encoded in low-level visual patterns but also in semantic content that can be described linguistically.

\paragraph{7. Absolute accuracies exceed those reported in the reference framework}

Across transformations shared with the Zeng et al. framework, the accuracies obtained here substantially exceed those reported for purely web-scraped dataset comparisons, by approximately 15--25 percentage points. While a direct comparison is not appropriate given the fundamental difference in dataset composition, this gap is itself informative. Zeng et al. analyse separability between web-scraped corpora that exhibit considerable internal diversity due to their heterogeneous photographic sources. In contrast, a generative model such as Stable Diffusion v1.5 imposes consistent learned priors over colour, texture, spatial layout, and scene composition across all its outputs, producing a statistical homogeneity that is absent in web-scraped corpora assembled from heterogeneous photographic sources. The substantially higher accuracies observed here, combined with the near-perfect separability retained even under the most aggressive compression, are therefore strongly consistent with the generative pipeline being the primary driver of separability in this work. While this remains an inference rather than a direct experimental finding, it is difficult to explain the pervasiveness of separability across every representational level without appeal to the systematic influence of the generative process.

% Across comparable transformations, the obtained accuracies exceed those reported by Zeng et al. by approximately 15–25 percentage points. This discrepancy is most likely attributable to dataset composition rather than methodological differences.

% The reference framework analyses separability between multiple web-scraped image datasets, which aggregate images from a wide range of photographic sources and therefore exhibit considerable internal diversity. In contrast, the dataset combination studied here includes Stable Diffusion v1.5 generated images, which are produced through a consistent generative pipeline that imposes shared priors over visual structure, texture statistics, and colour distributions.

% The presence of a generative dataset therefore introduces a class whose visual statistics are influenced by the generative process, potentially producing more pronounced dataset-level signatures than those observed in purely web-scraped corpora. Consequently, the results reported here should be interpreted in the context of this specific three-dataset configuration rather than as a direct replication of the reference-paper findings.

% Furthermore, the near-perfect separability observed across virtually every transformation — particularly between Re-LAION-5B and Stable Diffusion v1.5 — is strongly consistent with the generative pipeline being the primary driver of separability. A generative model imposes consistent learned priors over colour, texture, spatial layout, and scene composition across all its outputs, producing a statistical homogeneity that is absent in web-scraped corpora assembled from heterogeneous photographic sources. While this remains an inference rather than a direct experimental finding, the pervasiveness of separability across every representational level is difficult to explain without appeal to the systematic influence of the generative process.

\section{Demographic Attribute Separability}
\label{sec:demographic_attribute_separability}

To investigate whether demographic attributes are encoded in non-appearance visual
features, a series of probe experiments were conducted using transformations designed to
isolate distinct information channels: spatial structure, semantic content, body geometry,
and global colour statistics. These probes follow the same methodological framework
applied in the dataset classification experiments, substituting demographic labels for
dataset identity as the classification target. Crucially, the occlusion family of probes
directly tests whether demographic signal persists in the scene context once the person
region is removed, providing evidence about spurious correlations between demographic
attributes and non-person visual features. Results are reported as AUC, with chance level
at 0.500 for binary gender classification and 0.500 macro-averaged AUC for eight-class
skin tone classification.

To assess result stability, each probe was trained independently 
across four random seeds (0, 1, 2, 3); the best-performing seed 
result is reported for each probe and dataset combination. 
Variation across seeds was negligible for near-chance conditions 
on Re-LAION-5B, confirming the absence of recoverable signal 
rather than seed sensitivity.

MST categories 9 and 10 are absent from both datasets: no 
images in either corpus received these labels under the VGG-16 
MST10 Regression annotator. This outcome is consistent with the 
distributional analysis in Section~\ref{sec:demographic_distributions}, 
which shows near-zero Dark-bin proportions in both corpora. The 
annotator is capable of assigning these categories as confirmed during annotator evaluation, making it likely that MST 9 and 10 skin tones are extremely rare or entirely absent in the profession-conditioned image sets 
rather than systematically missed by the classifier. All probing 
experiments therefore operate over MST categories 1--8.

%-----------------------------------------------------------------------
\subsection{Gender Artifact Analysis}
\label{subsec:gender_artifact_analysis}
%-----------------------------------------------------------------------

\begin{table}[htbp]
	\centering
	\caption{Gender separability results per image transformation and dataset.
		All probes trained with seed 0 unless otherwise noted. AUC is the primary metric; chance level
		is 0.500 for binary classification.}
	\label{tab:gender_separability}
	\small
	\begin{threeparttable}
	\resizebox{\textwidth}{!}{%
	\begin{tabular}{llcc}
		\toprule
		\textbf{Transformation} & \textbf{Probe} &
		\textbf{Re-LAION-5B AUC} & \textbf{SD v1.5 AUC} \\
		\midrule
		Occlusion (full, no background)             & ResNet-50           & 0.499$^\dagger$ & 0.948 \\
		Occlusion (masked rectangle)                & ResNet-50           & 0.501$^\dagger$ & 0.747 \\
		Occlusion (masked rectangle, no background) & ResNet-50           & 0.499$^\dagger$ & 0.587$^\ddagger$ \\
		Occlusion (masked segmentation)             & ResNet-50           & 0.501$^\dagger$ & 0.852 \\
		Occlusion (masked segmentation, no bg.)     & ResNet-50           & 0.501$^\dagger$ & 0.834 \\
		Depth at pose keypoints                     & Logistic Regression & 0.559           & 0.592 \\
		Mean RGB                                    & Logistic Regression & 0.600           & 0.579 \\
		Object detection (restricted)               & Logistic Regression & 0.615           & 0.608 \\
		Object detection                            & Logistic Regression & 0.692           & 0.608 \\
		Pose keypoints                              & MLP        		  & 0.697           & 0.680 \\
		Base (unmodified)                           & ResNet-50           & 0.986           & 0.946 \\
		Captions									& SentenceT5          & 0.901			& 0.771 \\
		\bottomrule
	\end{tabular}}
	\begin{tablenotes}
		\footnotesize
		\item[$\dagger$] AUC at chance level; model collapsed to majority-class prediction.
		\item[$\ddagger$] Near-chance AUC; recall $\approx 1.0$ indicates near-complete majority-class prediction.
	\end{tablenotes}
	\end{threeparttable}
\end{table}

\paragraph{Unmodified images establish a strong separability upper bound}

The base ResNet-50 probe achieves AUC 0.986 on Re-LAION-5B and 0.946 on SD v1.5,
confirming that gender-correlated visual signal is encoded at high levels in both datasets
before any transformation is applied. The slightly higher value on Re-LAION-5B likely
reflects the greater photographic diversity of a web-scraped corpus, which produces a
wider variety of appearance cues correlated with gender labels, such as facial structure,
body morphology, hairstyle, and clothing. These results establish an upper bound against
which all reduced representations are compared and confirm that the annotation pipeline
produces labels sufficiently consistent to support downstream probing.

\paragraph{Occlusion probes reveal a fundamental asymmetry between datasets}

The occlusion family of probes constitutes the most consequential set of results in this
analysis, as it directly tests whether gender information is localised within the person
or distributed across the broader scene context.

% On Re-LAION-5B, every occlusion variant collapses performance to chance (AUC
% $\approx$ 0.500), regardless of whether masking uses a bounding rectangle or a tight
% segmentation contour, and regardless of whether the background is retained or removed.
% This uniform collapse establishes that gender-correlated signal in Re-LAION-5B resides
% exclusively within the visible appearance of the depicted person. No residual signal is
% recoverable from background objects, spatial layout, or any other contextual feature once
% the person region is occluded.

On Re-LAION-5B, all person-hidden occlusion variants collapse 
to chance (AUC $\approx$ 0.500), regardless of whether masking 
uses a bounding rectangle or a tight segmentation contour, and 
regardless of whether the background is retained or removed. 
This establishes that background context alone carries no 
recoverable gender signal in Re-LAION-5B. The full-person-no-background 
condition also yields chance performance (AUC 0.499), though this 
result is less straightforward to interpret: the probe collapses 
to majority-class prediction, which may reflect the increased 
difficulty of learning from person crops under the high 
photographic variability of web-scraped imagery, annotation 
dependency on contextual features, or a genuine absence of 
learnable person-level signal in isolation. The unambiguous 
finding is that background context carries no independent 
demographic signal in Re-LAION-5B.

The pattern for SD v1.5 is qualitatively different and constitutes the central finding of this analysis. Both the scene background and the person silhouette are substantial independent carriers of gender signal. When the person is masked with a rectangle and the background is retained, the probe achieves AUC 0.747 from background context alone. When the segmentation mask is applied without any background — leaving the classifier with only the person-shaped occluded region and no surrounding scene — performance rises to AUC 0.834, a 8.7 percentage-point increase over the rectangle-with-background condition. The silhouette outline itself therefore encodes gender signal beyond what background context alone provides, likely through the body shape and spatial proportions preserved by the tight contour, and is marginally the stronger individual signal source in isolation. When both silhouette and background are simultaneously available under segmentation masking, performance reaches 0.852, confirming that both sources contribute additively. The most consequential finding, however, is not which source is stronger in isolation, but that background context carries substantial recoverable gender signal at all — a pattern entirely absent in Re-LAION-5B, where every person-hidden condition collapses to chance regardless of whether the background is retained or removed. This directly replicates the contextual gender artifact identified by Meister et al. and demonstrates that the effect is substantially more pronounced in generatively produced images than in web-scraped corpora.

% The pattern for SD v1.5 is qualitatively different and constitutes the central finding of
% this analysis. The primary carrier of gender signal is the scene background: when the
% person is masked with a rectangle and the background is retained, the probe achieves
% AUC~0.747 from background context alone. However, the person-shaped silhouette is also
% independently informative. When the segmentation mask is applied without any background
% --- leaving the classifier with only the person-shaped occluded region and no surrounding
% scene --- performance rises to AUC~0.834, a 8.7 percentage-point increase over the
% rectangle-with-background condition. The silhouette outline itself therefore encodes
% gender signal beyond what background context alone provides, likely through the body
% shape and spatial proportions preserved by the tight contour. When both silhouette and
% background are simultaneously available under segmentation masking, performance reaches
% 0.852, confirming that both sources contribute additively.

% The strongest evidence for the background as the dominant signal source comes from
% the sharpest drop in the table: rectangle masking with background removal, which
% eliminates both the person region and the immediately surrounding context while
% retaining only the distal image periphery, reduces performance to AUC~0.587. Removing
% the background is therefore far more damaging to classification performance than
% any other single intervention, directly demonstrating that the scene context is the
% primary carrier of gender information in SD v1.5 images. This finding directly replicates
% the contextual gender artifact identified by Meister et al.\ and demonstrates that the
% effect is substantially more pronounced in generatively produced images than in
% web-scraped corpora.

\paragraph{Global colour statistics carry modest but consistent signal in both datasets}

Mean RGB classification achieves AUC 0.600 on Re-LAION-5B and 0.579 on SD v1.5. Both
values are in close agreement with the AUC of 0.580 reported by Meister et al.\ for mean
RGB gender probing on COCO, providing a useful external reference point for the magnitude
of this signal. Both datasets therefore exhibit gender-correlated photometric statistics
at a level consistent with a general-purpose image benchmark, indicating that systematic
differences in clothing colour or other photometric properties associated with gender
labels produce detectable signal even under compression to three global scalars. The
near-identical performance across datasets further suggests that this source of signal
reflects how occupational images are composed generally, rather than a characteristic
specific to either corpus.

\paragraph{Geometric probes reveal body configuration as an independent signal source}

Pose keypoints achieve AUC 0.697 on Re-LAION-5B and 0.680 on SD v1.5, indicating that
normalised skeletal geometry correlates measurably with gender labels in both datasets.
Because pose keypoints encode body proportions, posture, and limb arrangement without
conveying colour or texture information, this result demonstrates that systematic
differences in depicted body configurations contribute to gender separability
independently of photometric properties.

Depth values sampled at pose keypoint locations yield AUC 0.559 on Re-LAION-5B and 0.592
on SD v1.5 --- the lowest values recorded for any probe that produces above-chance
predictions. Although both exceed chance, the magnitude of separation is small, indicating
that 3D geometric structure at skeletal locations contributes only weak gender-correlated
signal. This positions depth at keypoints as a substantially less informative probe than
skeletal geometry alone, and suggests that the dominant spurious features in these
datasets are photometric and contextual in nature rather than geometric.

\paragraph{Semantic probes reveal scene composition and language as additional signal carriers}

Object detection probes show that scene composition carries gender-correlated information
beyond what body geometry alone explains. On Re-LAION-5B, the unrestricted object
detection probe achieves AUC 0.692, compared to 0.615 for the person-category-restricted
variant. This gap of 7.7 percentage points demonstrates that the objects appearing
alongside people in Re-LAION-5B reflect gender-stereotyped scene compositions
independently of the person's appearance. On SD v1.5, both variants converge at AUC
0.608, indicating that the generative model does not amplify this form of object-level
contextual stereotyping to the same degree as is observed for spatial background signal.

Caption-based classification using SentenceT5 embeddings achieves AUC 0.771 on SD v1.5,
demonstrating that gender-correlated signal survives compression to natural language.
Because the probe receives only textual input, this result indicates that the generative
pipeline systematically produces scene descriptions whose semantic content co-varies with
the gender of the depicted person: occupational contexts, object references, and activity
settings described in captions differ systematically across gender categories. The
persistence of this signal through a complete modality shift suggests that the stereotyping
identified through visual probes is coherent and cross-modal rather than confined to
low-level image statistics.

\paragraph{The two datasets exhibit qualitatively distinct spurious gender feature profiles}

% Considered jointly, the results reveal a structural divergence between the two datasets.
% In Re-LAION-5B, background context carries no independent gender signal: all 
% person-hidden occlusion variants collapse to chance regardless of whether the background 
% is retained or removed, and residual contributions from colour statistics, object 
% co-occurrence, and skeletal geometry are modest. The full-person-no-background condition 
% also collapses to chance, though this result is ambiguous and does not permit a strong 
% conclusion about person-level signal in isolation. In SD v1.5, gender information is 
% distributed across the full image: background context alone reproduces near-complete 
% classification performance, every region of the image contributes demographic signal, 
% and stereotyping is recoverable even from natural language descriptions. This profile 
% is consistent with the hypothesis that SD v1.5 amplifies the demographic priors of its 
% Re-LAION-based training corpus, projecting them not only onto the generated person but 
% onto the surrounding scene, the spatial layout, and the textual register of the 
% generated content.

Considered jointly, the results reveal a structural divergence between the two datasets. In Re-LAION-5B, background context carries no independent gender signal: all person-hidden occlusion variants collapse to chance regardless of whether the background is retained or removed, and residual contributions from colour statistics, object co-occurrence, and skeletal geometry are modest. The full-person-no-background condition also collapses to chance, though this result is ambiguous and does not permit a strong conclusion about person-level signal in isolation. In SD v1.5, gender information is uniquely distributed into scene context and background, background context alone carries substantial recoverable demographic signal, a pattern entirely absent in Re-LAION-5B. Additionally, stereotyping is recoverable even from natural language descriptions. While certain non-appearance representations such as mean RGB, object co-occurrence, and skeletal geometry carry comparable or weaker signal in SD v1.5 than in Re-LAION-5B, the defining characteristic of SD v1.5 is the presence of strong contextual background signal that has no counterpart in the web-scraped corpus. This profile is consistent with the hypothesis that SD v1.5 amplifies the demographic priors of its Re-LAION-based training corpus, projecting them not only onto the generated person but onto the surrounding scene, the spatial layout, and the textual register of the generated content.

%-----------------------------------------------------------------------
\subsection{Skin Tone Artifact Analysis}
%-----------------------------------------------------------------------

\begin{table}[htbp]
	\centering
	\caption{Skin tone separability results per image transformation and dataset
		(MST 1--8). All probes trained with seed 0 unless otherwise noted.
		Macro-averaged AUC is the primary metric; chance level is 0.500.}
	\label{tab:skintone_separability_overall}
	\small
	\begin{threeparttable}
	\resizebox{\textwidth}{!}{%
	\begin{tabular}{llcc}
		\toprule
		\textbf{Transformation} & \textbf{Probe} &
		\textbf{Re-LAION-5B AUC} & \textbf{SD v1.5 AUC} \\
		\midrule
		Occlusion (full, no background)             & ResNet-50           & 0.514$^\dagger$ & 0.899 \\
		Occlusion (masked rectangle)                & ResNet-50           & 0.514$^\dagger$ & 0.728 \\
		Occlusion (masked rectangle, no background) & ResNet-50           & 0.515$^\dagger$ & 0.660 \\
		Occlusion (masked segmentation)             & ResNet-50           & 0.514$^\dagger$ & 0.774 \\
		Occlusion (masked segmentation, no bg.)     & ResNet-50           & 0.515$^\dagger$ & 0.704 \\
		Mean RGB                                    & Logistic Regression & 0.563           & 0.543 \\
		Object detection (restricted)               & Logistic Regression & 0.591           & 0.626 \\
		Pose keypoints                              & MLP         & 0.596           & 0.666 \\
		Depth at pose keypoints                     & Logistic Regression & 0.598           & 0.655 \\
		Object detection                            & Logistic Regression & 0.611           & 0.663 \\
		Captions$^\ddagger$                         & SentenceT5          & 0.749           & 0.673 \\
		Base (unmodified)                           & ResNet-50           & 0.968           & 0.900 \\
		\bottomrule
	\end{tabular}}
 	\begin{tablenotes}
		\footnotesize
		\item[$\dagger$] Near-chance macro AUC; per-class AUCs at chance except MST8 (small-sample artefact).
		\item[$\ddagger$] Probe operates on BLIP Large captions encoded via SentenceT5-Base.
	\end{tablenotes}
	\end{threeparttable}
\end{table}

\begin{table}[htbp]
	\centering
	\caption{Per-class skin tone AUC per image transformation and dataset
		(MST 1--8). MST~9 and MST~10 are absent from both datasets and omitted.
		All probes trained with seed 0 unless otherwise noted.}
	\label{tab:skintone_separability_perclass}
	\small
	\begin{threeparttable}
	\resizebox{\textwidth}{!}{%
	\begin{tabular}{llcccccccccc}
		\toprule
		\textbf{Transf.} & \textbf{Dataset} &
		\textbf{MST1} & \textbf{MST2} & \textbf{MST3} & \textbf{MST4} &
		\textbf{MST5} & \textbf{MST6} & \textbf{MST7} & \textbf{MST8} &
		\textbf{MST9} & \textbf{MST10} \\
		\midrule
		Occ.\ (full., no bg)  & Re-LAION-5B & 0.499 & 0.501 & 0.501 & 0.499 & 0.501 & 0.499 & 0.499 & 0.614 & --- & --- \\
		Occ.\ (full., no bg)  & SD v1.5		& 0.969 & 0.936 & 0.881 & 0.778 & 0.783 & 0.888 & 0.963 & 0.996 & --- & --- \\
		Occ.\ (rect.)         & Re-LAION-5B & 0.499 & 0.499 & 0.499 & 0.499 & 0.501 & 0.501 & 0.499 & 0.614 & --- & --- \\
		Occ.\ (rect.)         & SD v1.5     & 0.837 & 0.777 & 0.715 & 0.594 & 0.625 & 0.685 & 0.741 & 0.853 & --- & --- \\
		Occ.\ (rect., no bg.) & Re-LAION-5B & 0.499 & 0.501 & 0.501 & 0.501 & 0.501 & 0.499 & 0.501 & 0.614 & --- & --- \\
		Occ.\ (rect., no bg.) & SD v1.5     & 0.783 & 0.679 & 0.631 & 0.538 & 0.586 & 0.595 & 0.618 & 0.846 & --- & --- \\
		Occ.\ (seg.)          & Re-LAION-5B & 0.499 & 0.499 & 0.499 & 0.501 & 0.499 & 0.501 & 0.499 & 0.614 & --- & --- \\
		Occ.\ (seg.)          & SD v1.5     & 0.876 & 0.821 & 0.751 & 0.625 & 0.645 & 0.738 & 0.815 & 0.921 & --- & --- \\
		Occ.\ (seg., no bg.)  & Re-LAION-5B & 0.501 & 0.501 & 0.499 & 0.501 & 0.501 & 0.499 & 0.501 & 0.614 & --- & --- \\
		Occ.\ (seg., no bg.)  & SD v1.5     & 0.817 & 0.719 & 0.669 & 0.581 & 0.596 & 0.655 & 0.712 & 0.883 & --- & --- \\
		Mean RGB              & Re-LAION-5B & 0.546 & 0.586 & 0.570 & 0.576 & 0.518 & 0.563 & 0.583 & 0.558 & --- & --- \\
		Mean RGB              & SD v1.5     & 0.585 & 0.594 & 0.525 & 0.520 & 0.529 & 0.538 & 0.551 & 0.499 & --- & --- \\
		Obj.\ det.\ (restr.)  & Re-LAION-5B & 0.682 & 0.596 & 0.578 & 0.550 & 0.540 & 0.548 & 0.577 & 0.653 & --- & --- \\
		Obj.\ det.\ (restr.)  & SD v1.5     & 0.691 & 0.648 & 0.630 & 0.538 & 0.579 & 0.599 & 0.627 & 0.692 & --- & --- \\
		Pose keypoints        & Re-LAION-5B & 0.691 & 0.633 & 0.576 & 0.565 & 0.554 & 0.543 & 0.573 & 0.630 & --- & --- \\
		Pose keypoints        & SD v1.5     & 0.770 & 0.668 & 0.640 & 0.539 & 0.588 & 0.610 & 0.638 & 0.875 & --- & --- \\
		Depth at pose         & Re-LAION-5B & 0.730 & 0.638 & 0.581 & 0.560 & 0.541 & 0.541 & 0.565 & 0.624 & --- & --- \\
		Depth at pose         & SD v1.5     & 0.756 & 0.657 & 0.630 & 0.536 & 0.583 & 0.602 & 0.615 & 0.859 & --- & --- \\
		Obj.\ det.            & Re-LAION-5B & 0.705 & 0.639 & 0.586 & 0.564 & 0.550 & 0.564 & 0.588 & 0.691 & --- & --- \\
		Obj.\ det.            & SD v1.5     & 0.735 & 0.683 & 0.649 & 0.566 & 0.589 & 0.637 & 0.672 & 0.775 & --- & --- \\
		Captions$^\ddagger$   & Re-LAION-5B & 0.784 & 0.732 & 0.719 & 0.760 & 0.734 & 0.751 & 0.745 & 0.769 & --- & --- \\
		Captions$^\ddagger$   & SD v1.5     & 0.772 & 0.688 & 0.653 & 0.582 & 0.594 & 0.653 & 0.692 & 0.747 & --- & --- \\
		Base                  & Re-LAION-5B & 0.976 & 0.976 & 0.969 & 0.961 & 0.948 & 0.959 & 0.973 & 0.980 & --- & --- \\
		Base                  & SD v1.5     & 0.966 & 0.935 & 0.883 & 0.779 & 0.783 & 0.895 & 0.964 & 0.996 & --- & --- \\
		\bottomrule
	\end{tabular}}
	\begin{tablenotes}
	\footnotesize
	\item[$\ddagger$] Probe operates on BLIP Large captions encoded via SentenceT5-Base.
	\end{tablenotes}
	\end{threeparttable}
\end{table}

\paragraph{Unmodified images establish a strong separability upper bound}

The base ResNet-50 probe achieves macro-AUC 0.968 on Re-LAION-5B and 0.900 on SD v1.5,
confirming that skin tone information is strongly encoded in person appearance within both
datasets. The per-class results are notably consistent across MST categories for
Re-LAION-5B (range 0.948--0.980), whereas SD v1.5 exhibits considerably wider variation
(0.779--0.996), with the lowest values in the mid-range tones MST4 and MST5. This
disparity indicates that the generative model renders skin tone signal less uniformly
across the tonal spectrum than is observed in real imagery, a tendency that becomes more
pronounced under reduced representations.

\paragraph{Occlusion probes reveal the same dataset asymmetry observed for gender,
	with per-class variation}

% All occlusion variants on Re-LAION-5B produce macro-AUC values at chance (0.514--0.515).
% The per-class breakdown confirms this is not an artefact of cancellation: individual
% AUCs for MST1 through MST7 are uniformly at chance (0.499--0.501) under every masking
% strategy. The exception is MST8, which registers AUC 0.614 consistently across all
% variants including those where the background is fully removed; given that this value
% is stable across masking strategies, it most plausibly reflects a small-sample artefact
% in this category rather than genuine contextual signal. The results establish that skin
% tone information in Re-LAION-5B is localised entirely within the visible person region,
% replicating the finding observed for gender in this dataset.

All occlusion variants on Re-LAION-5B produce macro-AUC values 
at chance (0.514--0.515). The per-class breakdown confirms this 
is not an artefact of cancellation: individual AUCs for MST1 
through MST7 are uniformly at chance (0.499--0.501) under every 
masking strategy. The exception is MST8, which registers AUC 
0.614 consistently across all variants including those where the 
background is fully removed; given that this value is stable 
across masking strategies, it most plausibly reflects a 
small-sample artefact in this category rather than genuine 
contextual signal. All person-hidden variants confirm that 
background context alone carries no recoverable skin tone signal 
in Re-LAION-5B. As with gender, the full-person-no-background 
condition also collapses to chance (0.514), though this result 
is ambiguous: it may reflect the photographic variability of 
web-scraped person crops, annotation dependency on contextual 
features, or a genuine absence of learnable person-level signal 
in isolation. The unambiguous finding is that background context 
carries no independent skin tone signal in Re-LAION-5B, 
replicating the pattern observed for gender.

On SD v1.5, occlusion probes retain substantial above-chance performance across all
masking strategies, with macro-AUC ranging from 0.660 to 0.899. The per-class results
reveal that this contextual signal is concentrated at the extremes of the tonal
distribution: MST1 and MST8 achieve the highest per-class AUCs across all variants
(e.g., 0.876 and 0.921 under segmentation masking), while mid-range tones MST4 and MST5
remain substantially lower. This indicates that the generative model encodes the strongest
contextual skin tone priors at the tonal extremes, where stereotypical scene associations
may be most pronounced in the training data.

The same occlusion gradient observed for gender is present here: segmentation masking
consistently outperforms rectangle masking, and retaining the background alongside person
occlusion consistently outperforms background removal. Both patterns confirm that the
immediately surrounding scene is the primary carrier of contextual skin tone signal in
SD v1.5.

\paragraph{Global colour statistics carry weaker photometric signal than for gender}

Mean RGB achieves macro-AUC 0.563 on Re-LAION-5B and 0.543 on SD v1.5, noticeably lower
than the corresponding gender values of 0.600 and 0.579. This reduction is consistent
with the expectation that global colour statistics carry stronger signal for an attribute
whose visual correlates include clothing and hairstyle than for one whose primary cue is a localised photometric property of skin regions: averaging colour across the entire image
dilutes the localised skin tone signal more severely than it dilutes the broader
gender-correlated colour variation. Per-class mean RGB AUCs are broadly distributed
across all MST categories for Re-LAION-5B (range 0.518--0.586), with no clear
concentration at tonal extremes, suggesting that colour-based skin tone stereotyping at
the global level is spread across categories, rather than category-specific.

\paragraph{Geometric probes produce comparable signal for skin tone and gender, with
	SD v1.5 showing stronger encoding}

Pose keypoints and depth at pose keypoints yield macro-AUC values of 0.596 and 0.598
respectively on Re-LAION-5B, and 0.666 and 0.655 on SD v1.5. Notably, these values are 
close to one another within each dataset, indicating that skeletal geometry and 3D keypoint depth carry roughly
equivalent amounts of skin tone information. The consistently higher SD v1.5 figures
across both probes reinforce the pattern that the generative model encodes demographic
attributes more pervasively into structural representations than the web-scraped corpus.

The per-class breakdown for pose keypoints reveals that, as with the occlusion probes,
the strongest signal in SD v1.5 is concentrated in MST1 and MST8 (AUC 0.770 and 0.875
respectively), while MST4 remains close to 0.539. This convergent finding across multiple
probe types suggests that the systematic scene stereotyping identified under occlusion
propagates through body configuration priors as well as background context.

\paragraph{Object detection probes show that scene composition encodes skin tone in
	both datasets}

Both object detection variants achieve moderate macro-AUC values across datasets:
the restricted probe yields 0.591 on Re-LAION-5B and 0.626 on SD v1.5, while the
unrestricted probe achieves 0.611 and 0.663 respectively. Unlike the gender case, where
the unrestricted probe on Re-LAION-5B substantially outperformed the restricted variant
(0.692 vs.\ 0.615), the gap here is modest on both datasets (approximately 2 percentage
points for Re-LAION-5B, 3.7 for SD v1.5), indicating that beyond-person object
co-occurrence contributes less additional skin tone signal than it does for gender. The
consistently higher SD v1.5 values, however, confirm that the generative model encodes
skin tone more strongly in scene-level object distributions than the web-scraped corpus,
consistent with the pattern across all non-appearance probes.

% \paragraph{Caption probing reveals a cross-dataset asymmetry absent from the visual
% 	probes}

% Caption-based classification produces a striking reversal of the pattern observed for
% all visual probes: Re-LAION-5B achieves macro-AUC 0.749, substantially higher than SD
% v1.5 at 0.673. This asymmetry --- where the real dataset exhibits stronger linguistic
% skin tone signal than the generated dataset --- directly contrasts with the visual
% occlusion results, where SD v1.5 substantially outperformed Re-LAION-5B.

% The Re-LAION-5B per-class caption AUCs are notably consistent across tones
% (range 0.719--0.784), indicating that skin tone stereotyping at the linguistic level is
% distributed broadly across the full tonal spectrum rather than concentrated at the
% extremes. For SD v1.5, per-class caption AUCs show the familiar concentration at MST1
% and MST8, while mid-range tones produce lower values, mirroring the pattern seen in the
% visual probes. Together, these results suggest that while SD v1.5 embeds skin tone priors
% more pervasively into visual scene statistics than Re-LAION-5B, the corresponding
% linguistic stereotyping in the captions of real images is no weaker --- and for mid-range
% tones, stronger --- than in the generated dataset. Skin tone stereotyping in web-scraped
% data may therefore be more evenly distributed across the tonal spectrum at the semantic
% level than the visual probe results alone would suggest.

\paragraph{Caption probing replicates the cross-dataset reversal observed for gender}
Caption-based classification replicates the asymmetry already identified in the gender caption results: Re-LAION-5B achieves macro-AUC 0.749, substantially higher than SD v1.5 at 0.673. As with gender captions, this inverts the pattern seen across all visual probes, where SD v1.5 consistently exhibited stronger demographic signal than Re-LAION-5B. The convergence of this reversal across both demographic attributes strengthens the interpretation that it reflects a structural difference in how demographic associations are encoded linguistically in real versus generated images, rather than an artefact specific to skin tone.

% The Re-LAION-5B per-class caption AUCs are notably consistent across tones (range 0.719--0.784), indicating that skin tone stereotyping at the linguistic level is distributed broadly across the full tonal spectrum rather than concentrated at the extremes. For SD v1.5, per-class caption AUCs show the familiar concentration at MST1 and MST8, while mid-range tones produce lower values, mirroring the pattern seen in the visual probes. Together, these results suggest that while SD v1.5 embeds skin tone priors more pervasively into visual scene statistics than Re-LAION-5B, the corresponding linguistic stereotyping in the captions of real images is no weaker than in the generated dataset. Skin tone stereotyping in web-scraped data may therefore be more evenly distributed across the tonal spectrum at the semantic level than the visual probe results alone would suggest.

The Re-LAION-5B per-class caption AUCs are notably consistent across tones (range 0.719--0.784), indicating that skin tone stereotyping at the linguistic level is distributed broadly across the full tonal spectrum rather than concentrated at the extremes. For SD v1.5, per-class caption AUCs show the familiar concentration at MST1 and MST8, while mid-range tones produce lower values, mirroring the pattern seen in the visual probes. Together, these results suggest that while SD v1.5 embeds skin tone priors more pervasively into visual scene statistics than Re-LAION-5B, the corresponding linguistic stereotyping in the captions of real images is no weaker — and for mid-range tones, stronger — than in the generated dataset. Skin tone stereotyping in web-scraped data may therefore be more evenly distributed across the tonal spectrum at the semantic level than the visual probe results alone would suggest.

\paragraph{Per-class results reveal systematic difficulty for mid-range skin tones
	across all probes}

Across all non-appearance probes and both datasets, a consistent ordering emerges: MST1
and MST8 achieve the highest per-class AUCs while MST4 and MST5 produce the lowest,
with the disparity most pronounced in SD v1.5. This pattern reflects the greater visual
similarity of mid-range tones to neighbouring categories, which reduces the distinctiveness
of any indirect correlates. It also raises a methodological concern: audit methods relying
solely on macro-averaged metrics may mask substantially weaker separability for the
categories that are simultaneously most visually ambiguous and most likely to be
underrepresented in training data, potentially understating bias for mid-range skin tones.

\paragraph{Skin tone and gender exhibit analogous but not identical spurious feature
	profiles}

% The skin tone results largely recapitulate the structural pattern observed for gender:
% Re-LAION-5B localises demographic signal within the person region, while SD v1.5
% distributes it throughout the image. The principal differences are quantitative rather
% than qualitative --- skin tone contextual AUCs under occlusion are somewhat lower than
% those for gender, mean RGB is a weaker probe for skin tone, and the caption asymmetry
% reverses direction --- but the fundamental finding holds for both attributes: SD v1.5
% systematically projects demographic priors into scene context, object composition, body
% geometry, and linguistic register in a manner that is not observed in the web-scraped
% corpus.

% The skin tone results largely recapitulate the structural pattern observed for gender:
% background context carries no independent skin tone signal in Re-LAION-5B, while 
% SD v1.5 distributes demographic signal throughout the image. The principal differences 
% are quantitative rather than qualitative, skin tone contextual AUCs under occlusion 
% are somewhat lower than those for gender, mean RGB is a weaker probe for skin tone, 
% and the caption asymmetry reverses direction, but the fundamental finding holds for 
% both attributes: SD v1.5 systematically projects demographic priors into scene context, 
% object composition, body geometry, and linguistic register in a manner that is not 
% observed in the web-scraped corpus.

The skin tone results largely recapitulate the structural pattern observed for gender: background context carries no independent skin tone signal in Re-LAION-5B, while SD v1.5 uniquely distributes demographic signal into scene context and background. The principal differences are quantitative rather than qualitative, skin tone contextual AUCs under occlusion are somewhat lower than those for gender, mean RGB is a weaker probe for skin tone, and the caption asymmetry reverses direction, but the fundamental finding holds for both attributes: SD v1.5 systematically projects demographic priors into scene context, object composition, body geometry, and linguistic register in a manner that is not observed in the web-scraped corpus.

%-----------------------------------------------------------------------
\subsection{Connections to Dataset Classification Results}
%-----------------------------------------------------------------------

Several structural parallels link the demographic probe experiments 
to the earlier dataset classification findings, and together the two 
sets of results provide a mutually reinforcing account of how 
statistical bias is distributed within and between the datasets. 
% Given the uniformly high overall accuracies observed across all 
% transformations, near-perfect separability of each individual class 
% can be inferred without requiring explicit per-class reporting: at 
% 95\%+ overall accuracy under a balanced three-class prior, no single 
% class can be pulling performance substantially below chance while the 
% others achieve high accuracy.

The most direct parallel is the persistence of signal across 
representation levels. In the dataset classification experiments, 
separability remained above chance across every transformation 
tested, from global colour statistics at 64.19\% through to natural 
language at 90.99\%, indicating that dataset identity is encoded in 
multiple correlated visual cues simultaneously. The demographic 
probing results recapitulate this structure: neither gender nor skin 
tone is recoverable from a single isolated representation, but both 
remain above chance across photometric (mean RGB), geometric (pose 
keypoints, depth at keypoints), semantic (object detection, 
captions), and contextual (occlusion) probes. The implication in 
both cases is the same: the statistical regularities in question are 
not artefacts of a single image property but are embedded across the full representational spectrum.

The sharpest connection between the two analyses concerns the role 
of Stable Diffusion v1.5. In the dataset classification task, the 
presence of SDv1.5 as one of the three classes was proposed as the 
primary explanation for why classification accuracies substantially 
exceed those reported by Zeng et al.\ for purely web-scraped 
datasets: the generative pipeline imposes consistent visual priors 
that make SDv1.5 images far more internally homogeneous than 
Re-LAION-5B or COCO, and therefore far easier to identify as a 
class. The occlusion results in the demographic probing experiments 
provide a mechanistic account of this homogeneity. SDv1.5 encodes 
gender at AUC 0.948 in the unoccluded person and AUC 0.747 from 
background context alone; the corresponding skin tone values are 
0.900 and 0.728. These figures confirm that the generative model 
has learned coherent joint priors over person appearance and scene 
context that are stable enough across 200 professions to be 
exploited by a linear classifier. Re-LAION-5B shows no such contextual signal under any person-hidden masking condition, indicating that its scene context carries no independent demographic signal. The full-person-no-background condition also collapses to chance, though this result is 
ambiguous and does not permit a definitive conclusion about whether person-level signal 
is independently sufficient in real photographs. The dataset classification experiments cannot reveal which class is responsible for the 
high separability; the demographic probing experiments establish 
that the answer is predominantly SDv1.5.

% Re-LAION-5B shows no such 
% contextual signal under any masking condition, indicating that its 
% demographic information is confined to person pixels and that its 
% scene context is essentially uninformative about the gender or skin 
% tone of the depicted individual. 
% The dataset classification experiments cannot reveal which class is responsible for the 
% high separability; the demographic probing experiments establish 
% that the answer is predominantly SDv1.5.

The comparison between the two caption-based results further 
sharpens this picture. Dataset identity survived compression to 
natural language at 90.99\% classification accuracy, a figure 
that places caption-based classification at the same tier as 
semantic segmentation and above object detection in the 
discrimination hierarchy. Gender persists in captions at AUC 
0.771 for SDv1.5 and 0.901 for Re-LAION-5B; skin tone at AUC 
0.673 and 0.749 respectively. The reversal between datasets in 
the caption domain, where Re-LAION-5B shows stronger 
demographic separability than SDv1.5, in contrast to every 
visual probe, is the one result that does not replicate the 
pattern seen in the classification task. It suggests that 
while SDv1.5's generative priors produce stronger scene-level 
visual stereotyping, the real-image corpus contains richer 
linguistically recoverable demographic associations, likely 
because web-scraped images are drawn from a broader semantic 
space of human activity where occupational, environmental, and 
social contexts co-vary more systematically with demographic 
attributes at the description level. Taken together, the 
caption findings across both experiments confirm that 
stereotyping in these datasets is not confined to pixel 
statistics or geometric structure but extends into the 
semantic content of the imagery, recoverable from text alone.

Finally, the depth and edge classification results connect to 
the novel Depth-at-Pose-Keypoints representation introduced 
in the demographic probing pipeline. Depth achieves 96.78\% 
dataset classification accuracy and Canny edge 93.62\%, 
confirming that structural geometry carries strong 
dataset-discriminating signal. At the demographic level, depth 
at pose keypoints yields AUC 0.559 for gender and 0.598 for 
skin tone on Re-LAION-5B, and 0.592 and 0.655 on SDv1.5.
The contrast between the high dataset classification 
accuracy of full depth maps and the relatively weak demographic 
signal from sparse keypoint depth sampling suggests that the 
structural differences captured in full depth maps are 
primarily compositional, consisting of camera distance, scene depth 
layering and background structure rather than body-geometric, 
and that this compositional signal does not substantially 
co-vary with demographic attributes at the individual level. 
The sparse depth-at-keypoints representation, by isolating 
only the 3D configuration of the skeletal joints, extracts 
the demographically relevant portion of the structural signal 
while discarding the scene-level compositional information 
that drives the high classification accuracy.

%% ============================================================
%% Section 4: Debiasing Evaluation
%% ============================================================
\section{Debiasing Evaluation}
\label{sec:debiasing_evaluation}

The debiasing evaluation compares the demographic distribution and spurious feature
profile of the 10,000-image debiased doctor corpus against the full Stable Diffusion v1.5
generation results reported in Sections~\ref{sec:demographic_distributions}
and~\ref{sec:demographic_attribute_separability}. Rather than constructing a separate doctor-specific
baseline --- which would be expected to replicate or amplify the patterns already
identified across the 200,000-image standard corpus --- the evaluation uses the
aggregate SDv1.5 findings as the reference point. This comparison directly assesses
whether the combined ITI-GEN, ControlNet, and Chamfer colour alignment pipeline
produces a measurable departure from the demographic and spurious feature profile
characteristic of standard Stable Diffusion v1.5 generation.

An important methodological distinction applies throughout this section. Demographic
labels for the debiased corpus are not derived from automated classifiers but are assigned
directly from the generation configuration: each image is labelled with the ITI-GEN gender
token and MST skin tone token used to produce it. This eliminates annotation uncertainty as
a confound in the evaluation and ensures that the probing classifiers measure the degree
to which low-level image representations encode the demographic conditioning signal rather
than classifier-predicted attributes. For gender, high probing accuracy on person-level
representations is therefore an expected consequence of successful ITI-GEN conditioning
rather than evidence of spurious signal. For skin tone, the 10-class problem is
substantially harder than the binary gender case, and person-level probe values should be
interpreted relative to the chance level of 0.500 for macro-averaged AUC rather than in
direct comparison with the binary gender results.

\subsection{Demographic Distribution Shift}
\label{subsec:debias_distribution}

The debiased generation pipeline was configured to produce images uniformly distributed
across gender and skin tone categories. ITI-GEN embeddings were applied for male and
female gender categories in equal proportions across the 10,000-image corpus, with MST
skin tone conditioning applied uniformly across all ten MST levels (MST~1--10), yielding
a target distribution of approximately 500 images per gender--skin tone combination. After
structural validation using the same acceptance criteria as the standard generation
pipeline, the effective composition of the debiased corpus reflects this intended uniform
distribution, subject to profession-specific rejection rates.

By construction, the debiased corpus achieves a gender distribution of approximately 50\%
female and 50\% male. This stands in direct contrast to the strongly male-skewed output of
the standard Stable Diffusion v1.5 pipeline, which produced 29.02\% female images across
all 200 professions, a figure consistent with the amplification of training data gender
imbalance documented in Section~\ref{sec:demographic_distributions}.

The skin tone distribution of the debiased corpus is equally balanced by construction
across MST~1--10, in contrast to the pronounced light-bin skew of the standard pipeline,
where 76.99\% of generated images fell in the Mid bin and Dark-bin representation
approached zero across most professions. The debiasing pipeline therefore produces a
distribution that would require radical change to the model's learned prior to arise
naturally from standard generation.

Figure~\ref{fig:debias_sample_grid} shows a representative sample of debiased generated
images at a fixed seed, with rows corresponding to female and male gender conditioning
and columns corresponding to MST skin tone levels 1 through 10.

\begin{figure}[htbp]
	\centering
	\includegraphics[width=\textwidth]{debiased_grid.png}
	\caption[Representative debiased generation samples at a fixed seed.]{Representative
		debiased generation samples at a fixed seed. Rows correspond to female (top) and
		male (bottom) ITI-GEN gender conditioning; columns correspond to MST skin tone
		levels 1 through 10. Skin tone transitions progressively from lighter values at
		MST~1 to darker values at MST~10, demonstrating that the ten-point conditioning
		axis produces a visually coherent and continuous tonal gradient. Gender
		presentation is clearly differentiated between rows through facial features, body
		morphology, and clothing style. The ControlNet OpenPose conditioning enforces a
		consistent standing, full-body, centred pose across all twenty images, and the
		Chamfer colour alignment is reflected in the broadly consistent warm tonal palette
		maintained throughout the grid. The background exhibits some variation across
		columns, reflecting the residual scene-level variation that pose conditioning does
		not fully constrain, consistent with the background-level gender signal identified
		in Section~\ref{subsec:debias_probing}.}
	\label{fig:debias_sample_grid}
\end{figure}

\subsection{Probing Accuracy on the Debiased Corpus}
\label{subsec:debias_probing}

To assess whether the debiasing interventions reduced not only the aggregate demographic
imbalance but also the learnability of demographic signal from low-level image
representations, the same probing classifiers applied to the standard SDv1.5 corpus were
applied to the debiased corpus under identical experimental conditions.
Tables~\ref{tab:debias_gender_probing} and~\ref{tab:debias_skintone_probing} report gender
and skin tone probing AUC for the standard SDv1.5 corpus and the debiased corpus across
all evaluated representations.

\begin{table}[htbp]
	\centering
	\caption{Gender probing AUC on the standard Stable Diffusion v1.5 corpus and the
		debiased corpus. Chance level is 0.500 for binary classification. Labels for the
		debiased corpus are derived from the generation configuration rather than
		automated annotation. $\dagger$~Standard SDv1.5 value reflects near-complete
		majority-class collapse; see Section~\ref{subsec:gender_artifact_analysis} for detail.}
	\label{tab:debias_gender_probing}
	\begin{tabular}{llcc}
		\toprule
		\textbf{Representation} & \textbf{Probe} & \textbf{SDv1.5} & \textbf{Debiased} \\
		\midrule
		Base (unmodified)                          & ResNet-50           & 0.946 & $\approx$1.000 \\
		Occlusion (full, no background)            & ResNet-50           & 0.948 & 0.998 \\
		Occlusion (masked segmentation)            & ResNet-50           & 0.852 & 0.999 \\
		Occlusion (masked segmentation, no bg.)    & ResNet-50           & 0.834 & 0.985 \\
		Occlusion (masked rectangle)               & ResNet-50           & 0.747 & 0.995 \\
		Occlusion (masked rectangle, no bg.)       & ResNet-50           & 0.587$^\dagger$ & 0.701 \\
		\midrule
		Mean RGB                                   & Logistic Regression & 0.579 & 0.507 \\
		Object detection                           & Logistic Regression & 0.608 & 0.449 \\
		Object detection (restricted)              & Logistic Regression & 0.608 & 0.457 \\
		Depth at pose keypoints                    & Logistic Regression & 0.592 & 0.454 \\
		Pose keypoints                             & MLP                 & 0.680 & 0.475 \\
		Captions                                   & SentenceT5          & 0.771 & 0.530 \\
		\bottomrule
	\end{tabular}
\end{table}

\begin{table}[htbp]
	\centering
	\caption{Skin tone probing macro-averaged AUC on the standard Stable
		Diffusion v1.5 corpus and the debiased corpus. The SDv1.5 column reports
		8-class macro-averaged AUC (MST~1--8, as MST~9 and MST~10 are absent from
		the standard corpus); the debiased column reports 10-class macro-averaged AUC
		(MST~1--10, from the generation configuration). Chance level is 0.500 in both
		cases. Labels for the debiased corpus are derived from the generation
		configuration rather than automated annotation.}
	\label{tab:debias_skintone_probing}
	\begin{tabular}{llcc}
		\toprule
		\textbf{Representation} & \textbf{Probe} & \textbf{SDv1.5} & \textbf{Debiased} \\
		\midrule
		Base (unmodified)                          & ResNet-50           & 0.900 & 0.861 \\
		Occlusion (full, no background)            & ResNet-50           & 0.899 & 0.823 \\
		Occlusion (masked segmentation)            & ResNet-50           & 0.774 & 0.729 \\
		Occlusion (masked segmentation, no bg.)    & ResNet-50           & 0.704 & 0.663 \\
		Occlusion (masked rectangle)               & ResNet-50           & 0.728 & 0.632 \\
		Occlusion (masked rectangle, no bg.)       & ResNet-50           & 0.660 & 0.543 \\
		\midrule
		Mean RGB                                   & Logistic Regression & 0.543 & 0.482 \\
		Object detection                           & Logistic Regression & 0.663 & 0.439 \\
		Object detection (restricted)              & Logistic Regression & 0.626 & 0.461 \\
		Depth at pose keypoints                    & Logistic Regression & 0.655 & 0.443 \\
		Pose keypoints                             & MLP                 & 0.666 & 0.483 \\
		Captions                                   & SentenceT5          & 0.673 & 0.490 \\
		\bottomrule
	\end{tabular}
\end{table}

\subsubsection{Gender Probing Results}

Person-level representations confirm successful demographic conditioning. The unmodified
image probe achieves AUC $\approx$1.000 on the debiased corpus, and person-centric
occlusion variants are uniformly high: 0.998 for full person appearance without
background (FullNoBg), 0.999 for the segmentation-masked person region with background
retained (MaskSegm), and 0.985 for the silhouette alone without background
(MaskSegmNoBg). These values exceed the corresponding standard SDv1.5 figures by
substantial margins. As noted above, this is expected: ITI-GEN directly conditions the
person's visual appearance and labels are assigned from the conditioning token rather than
a classifier. High person-level separability confirms that the conditioning pipeline
successfully produced perceptually gendered images and that the demographic signal is
encoded at the level of the depicted person, precisely as intended.

Low-level spurious signal is substantially reduced across all non-person representations. Mean RGB classification falls from 0.579 on the standard SDv1.5 corpus to 0.507 on the debiased
corpus, effectively at the binary chance level of 0.500. Object detection probes fall to
0.449 and 0.457 for the unrestricted and restricted vocabularies respectively, both at or
near chance, indicating that the debiased corpus exhibits no systematic gender-correlated
object co-occurrence patterns. Depth at pose keypoints collapses to 0.454 and pose
keypoints to 0.475 (mean across five seeds), both within the range expected under a null
hypothesis of no demographic signal in geometric representations. These reductions stand
in direct contrast to the standard SDv1.5 results, where all of these probes achieved AUC
values between 0.579 and 0.680.

Caption-based probing falls from 0.771 on the standard SDv1.5 corpus to 0.530 on the
debiased corpus, a reduction of 24.1 percentage points. The near-halving of this signal
indicates that the debiasing pipeline substantially disrupts the generation of
gender-stereotyped semantic content, even though caption generation is not directly
conditioned by any of the debiasing interventions. This provides evidence that the
debiasing effect propagates into the semantic register of generated content.

The background-only probe (MaskRect) achieves AUC 0.995 on the debiased corpus, indicating that strong background-level signal persists. This is consistent with ITI-GEN's conditioning of scene content and the ControlNet pose constraint together influencing global composition in ways that correlate with the gender conditioning token. The standard SDv1.5 value of 0.747 for this condition provides the relevant comparison baseline. The persistence of near-person-level AUC in the background-only condition confirms that background-level demographic conditioning remains substantially entangled with the generation process, and represents the primary avenue for further improvement.

\subsubsection{Skin Tone Probing Results}

The skin tone probing results follow a consistent pattern of reduction across all
representations, with every probe falling relative to the standard SDv1.5 corpus. Unlike
the gender case, person-level probes decrease rather than increase in the debiased corpus:
the unmodified base probe falls from 0.900 to 0.861, and the full person appearance probe
(FullNoBg) falls from 0.899 to 0.823. This reduction reflects the substantially greater
difficulty of 10-class skin tone discrimination under a balanced uniform distribution
across MST~1--10 compared to the standard corpus, where the model's strong light-skewed
prior produces a concentrated distribution that is easier to learn. The person-level
values of 0.823--0.861 remain substantially above the 10-class chance level of 0.500,
confirming that the ITI-GEN conditioning successfully encodes skin tone signal at the
person level across the full ten-point scale.

The most direct evidence for the effectiveness of the Chamfer colour alignment is the
mean RGB probe, which falls from 0.543 on the standard SDv1.5 corpus to 0.482 on the
debiased corpus --- the skin tone probe closest to the 10-class chance level of
0.500. Since mean RGB compresses each image to its global average colour per channel,
a value at or near chance confirms that the debiased corpus exhibits no systematic
skin-tone-correlated colour statistics at the global level, directly demonstrating that
the Chamfer colour alignment disrupts the photometric correlations that enable colour
statistics to carry demographic signal.

All remaining non-person probes collapse to at or near chance. Object detection falls from 0.663 to 0.439 (unrestricted) and from 0.626 to 0.461 (restricted), both at or near chance, indicating that the debiased corpus exhibits no systematic skin-tone-correlated object co-occurrence patterns. Depth at pose keypoints falls from 0.655 to 0.443 and pose keypoints from 0.666 to 0.483, both at or near chance, consistent with no recoverable skin tone signal in geometric representations. Caption probing falls from 0.673 to 0.490, approaching chance, indicating that the semantic content of generated descriptions is no longer systematically correlated with the skin tone conditioning token.

The occlusion probes show consistent reductions across all variants. MaskRectNoBg falls from 0.660 to 0.543, close to chance, while MaskRect falls from 0.728 to 0.632, indicating a partial but incomplete reduction of background-level skin tone signal. The retention of above-chance signal in background-only probes suggests that the scene composition generated by the model retains some skin-tone-correlated properties even under pose conditioning, consistent with the background entanglement observed in the gender analysis.

% \subsubsection{Summary}

% Taken together, the gender and skin tone probing results establish a coherent and
% consistent pattern. The debiasing pipeline successfully conditions demographic attributes
% at the person level while substantially eliminating the low-level spurious correlations,
% geometric signal, and semantic stereotyping that characterise the standard Stable
% Diffusion v1.5 generation pipeline. The near-chance mean RGB skin tone probe provides
% the clearest mechanistic confirmation that the Chamfer colour alignment addresses
% its intended target. Residual background-level signal persists in both the gender and skin
% tone analyses, indicating that background-level demographic conditioning remains partially
% entangled with the generation process even under explicit pose and appearance control, and
% represents the primary avenue for further improvement.

%% ============================================================
%% Section 6: Limitations
%% ============================================================
\section{Limitations}
\label{sec:limitations}

Several limitations qualify the interpretation of the 
results presented in this chapter.

\paragraph{Scope of the Re-LAION-5B sample.}
The analysis is conducted on four of 128 shards of 
Re-LAION-5B, representing approximately 3.1\% 
of the full dataset. Findings are therefore estimates 
of corpus-level properties rather than exact 
population parameters, and the presence of unknown 
crawl-time biases cannot be excluded.

\paragraph{Annotation as measurement instrument.}
Demographic labels are produced by automated 
classifiers rather than ground-truth human annotation. 
Gender predictions are binary and do not capture 
non-binary gender expression. Skin tone estimates 
are sensitive to imaging conditions and lighting 
variability inherent in web-scraped data. Observed 
demographic distributions should therefore be 
interpreted as classifier outputs rather than 
definitive ground truth labels.

\paragraph{Dataset classification and causation.}
Above-chance classification accuracy confirms that 
a transformation retains dataset-discriminating 
signal, but does not establish which specific 
properties of each dataset produce that signal, 
nor does it constitute a causal claim about how 
those properties arose. The per-class accuracy 
attribution localises signal to a dataset but 
cannot decompose it further without additional 
controlled experiments.

\paragraph{Debiasing evaluation scope.}
The debiasing evaluation is restricted to a single profession and a single combination of inference-time interventions. It does not evaluate weight editing or fine-tuning approaches, and results may not generalise to other professions, prompts, or model architectures. The probing results demonstrate that low-level spurious correlations are substantially reduced at the representation level for the doctor profession, but whether this reduction generalises across the full 200-profession set or persists under different prompt formulations remains untested. The residual background-level signal identified in both the gender and skin tone analyses indicates that background-level demographic conditioning remains partially entangled with the generation process, and constitutes the primary open question for future work on this pipeline.

\paragraph{Synthetic image annotation reliability.}
Demographic classifiers and face detectors are 
trained on real photographs. Their reliability on 
synthetic images produced by latent diffusion 
models has not been formally validated in this 
work, and the presence of generative artefacts 
may affect annotation accuracy in ways that 
differ systematically between the real and 
synthetic corpora.


%% ============================================================
%% Section 7: Overall Conclusions
%% ============================================================
\section{Synthesis of Evaluation Findings}
\label{sec:overall_conclusions}

The dataset classification experiments establish that Re-LAION-5B, Stable Diffusion v1.5
generated images, and COCO are distinguishable with high accuracy across every level of
visual representation tested, from unmodified images at 99.42\% down to mean RGB
compression at 64.19\%. The persistence of above-chance separability even under the most
aggressive information reduction demonstrates that dataset identity is encoded in multiple
correlated visual cues simultaneously rather than in any single isolated artefact. The VAE
experiment confirms that this separability is not a consequence of low-level acquisition
artefacts, as reconstruction through the Stable Diffusion latent bottleneck reduces
accuracy by only 0.27 percentage points.

The demographic distribution analysis reveals that Re-LAION-5B exhibits gender distributions that diverge from real-world labour force composition in a baseline-dependent manner, and that Stable Diffusion v1.5 does not merely reflect the imbalances present in its training corpus but actively amplifies them across nearly every occupational group regardless of baseline. Gender amplification is confirmed in 178 of 200 professions, with the female proportion falling from 41.5\% in Re-LAION-5B to 29.02\% in Stable Diffusion v1.5. The skin tone analysis reveals a near-complete collapse of dark-toned representation in generative outputs: the Dark-bin proportion falls from 0.81\% in Re-LAION-5B to 0.14\% in Stable Diffusion v1.5, an 83\% reduction, while the Light-bin proportion almost doubles from 12.45\% to 22.87\%. Skin tone amplification toward lighter tones is confirmed in 175 of 200 professions.

The demographic probing experiments reveal qualitatively distinct spurious feature 
profiles between the two datasets. In Re-LAION-5B, background context carries no independent demographic signal: every person-hidden occlusion variant collapses to chance regardless of whether background is retained, and residual signal from colour statistics, object co-occurrence, and skeletal geometry is modest. The full-person-no-background condition also collapses to chance, though this result is ambiguous and does not support a strong conclusion about person-level signal in isolation. In Stable Diffusion v1.5, demographic information is uniquely distributed into scene context and background. It is recoverable from background context alone, from the person silhouette independently of background, and persists into natural language descriptions of generated content. This cross-modal persistence demonstrates that the stereotyping embedded in the generative model is coherent and pervasive rather than confined to low-level pixel statistics.

The debiasing evaluation demonstrates that the combined ITI-GEN, ControlNet, and Chamfer
colour alignment pipeline produces a measurable and substantive departure from the
demographic and spurious feature profile of the standard generation pipeline. Every
non-person probe collapses to at or near chance for both gender and skin tone, confirming
that the debiasing interventions substantially reduce the low-level spurious correlations, object-level stereotyping, and semantic stereotyping that characterise standard Stable Diffusion v1.5 generation. Residual background-level signal persists in both analyses, indicating that background-level demographic conditioning remains partially entangled with the generation process even under explicit pose and appearance control.